% Job posting: https://freehire.me/jobs/machine-learning-researcher-pushing-ai-frontiers-in-banking-rbc-3j5ep3ds
% =====================================================================
%  CV - Nishal Stanislaus Silva, PhD
%  Target: RBC Borealis (RBC's AI research institute) - Machine Learning Researcher
%  Variant: V2 (AI-Research) - publication-weighted, 3 pages
% =====================================================================

\documentclass[11pt]{article}
\usepackage[left=0.7in,top=0.7in,right=0.7in,bottom=0.7in]{geometry}
\usepackage{enumitem}
\usepackage[hidelinks]{hyperref}
\usepackage{titlesec}
\usepackage{array}
\usepackage{setspace}
\usepackage{needspace}
\usepackage[T1]{fontenc}
\usepackage{newtxtext,newtxmath}
\ifdefined\pdfgentounicode
  \input{glyphtounicode}
  \pdfgentounicode=1
\fi
\setstretch{1.04}
\pagenumbering{gobble}
\titleformat{\section}{\Large\scshape}{}{4em}{}[\vspace{-0.7em}\rule{\linewidth}{0.4pt}\vspace{-0.4em}]
\titlespacing*{\section}{0pt}{1em}{0.4em}
\newenvironment{cvrole}[5]{%
  \par\noindent\textbf{\large #1} | \textit{\small #2, #3}\hfill \textit{\small #4 - #5}\par
  \begin{itemize}[leftmargin=2em, itemsep=-0.3em, topsep=-0.5em]%
}{%
  \end{itemize}\par\noindent\mbox{}\par\vspace{0.1em}%
}
\newcommand{\cvName}{Nishal Stanislaus Silva}
\newcommand{\cvEmail}{nishal.silva@hotmail.com}
\newcommand{\cvPhone}{+1 613 200 2030}
\newcommand{\cvCity}{Toronto, ON}
\newcommand{\cvSite}{https://nishal.xyz}
\newcommand{\cvLinkedIn}{https://www.linkedin.com/in/nishal-silva}
\newcommand{\cvGitHub}{https://github.com/ns2max}
\newcommand{\cvStatus}{Canadian Permanent Resident, Eligible to work in Canada without sponsorship}
\newcommand{\cvTagline}{Research Scientist | Real-Time Multimodal ML, Benchmarks \& Open Datasets}

\begin{document}
\pagestyle{empty}
\begin{center}
    {\LARGE \scshape \textbf{\cvName}}{\large \textit{, PhD}}\\[1pt]
    {\small \cvTagline}\\[1pt]
    {\footnotesize\textit{\cvStatus}}\\[1pt]
    {\small
    \href{mailto:\cvEmail}{\cvEmail} \,\,|\,\, \cvPhone \,\,|\,\, \cvCity
    \,\,|\,\, \href{\cvSite}{nishal.xyz} \,\,|\,\, \href{\cvLinkedIn}{linkedin.com/in/nishal-silva}
    \,\,|\,\, \href{\cvGitHub}{github.com/ns2max}}
\end{center}

\section*{Professional Summary}

Machine learning researcher with a PhD (University of Trento, 2025) and 10+ years across industry and academia, centred on time-series and sequence modelling: RNN/LSTM/GRU architectures, benchmark and dataset design, and evaluation methodology that separates a real gain from a better-looking metric. I translate research into production-grade solutions: in my postdoc I owned both the science and the pipeline for a novel real-time system, publishing a detector that runs at F1 0.76 in 14 ms on a Raspberry Pi 4, 74$\times$ faster than the DTW baseline it replaced (JAES 2026); earlier, at MAS Holdings, I shipped research prototypes to the factory floor with CI/CD and code review. I prototype generative models for time series and signals (VAE- and diffusion-based synthetic data for limited-label regimes) and work routinely with large datasets: four open-access benchmark corpora (9,800+ recordings, 70+ participants) and high-frequency IoT ETL across three continents. My financial-anomaly-detection work at Forestpin covered forensic analytics, compliance monitoring and production scoring services. My record includes 10+ peer-reviewed papers, journal and program-committee review, and supervision of two BSc thesis students.

\section*{Core Competencies}

\textbf{Time-Series \& Sequence Modelling:} deep sequence models (RNN/LSTM/GRU, CNN), frozen-protocol benchmark and ablation design, evaluation under distribution shift, error analysis, statistical and probabilistic modelling, anomaly detection on messy structured and streaming data, DTW and classical baselines, uncertainty analysis\\
\textbf{Research to Production:} framing open-ended problems and delivering end-to-end systems with limited guidance, adapting techniques to latency/cost/reliability budgets, partnering with engineering and product to bring models to production, reproducible Python beyond notebooks (pytest, CI/CD), latency and compute characterization\\
\textbf{Generative Modelling \& Data:} VAE- and diffusion-based synthetic-data prototypes for time series and signals, rule-based variation generators, working with large datasets, four open-access benchmark datasets, high-frequency IoT ETL at manufacturing scale, multimodal sensor fusion (audio, IMU, EEG/BCI, XR telemetry)\\
\textbf{Publication \& Service:} 10+ peer-reviewed publications, reviewer for IEEE Access and the Journal of the National Science Foundation of Sri Lanka, program committee for IEEE IS\textsuperscript{2} and IEEE I3DA, supervision of two BSc thesis students, group-wide technical training and CI/CD practice at MAS

\section*{Technical Stack}

\textbf{Languages:} Python, C++17, C, MATLAB, SQL, JavaScript/TypeScript, Shell\\
\textbf{ML \& AI:} PyTorch, TensorFlow, Keras, scikit-learn, NumPy, Pandas, SciPy; RNN/LSTM/GRU, CNN, DTW baselines, anomaly detection, ablation and uncertainty analysis, VAE and diffusion prototypes; familiar with JAX, HuggingFace, ONNX, and LLM-assisted tooling (Claude Code)\\
\textbf{Signal \& Time-Series Processing:} DSP, STFT, MFCC, S-transform, CQT, onset and beat tracking, resampling and alignment, Librosa, SciPy, FFTW, music information retrieval\\
\textbf{Data \& Cloud:} AWS (EC2, S3, ECS, MWAA, RDS), Snowflake, SQL, ETL pipelines, NumPy, Pandas\\
\textbf{Dev \& MLOps:} Docker, Airflow, Jenkins, Git, Linux/UNIX, pytest, CI/CD, REST APIs

\section*{Education}
\textbf{Ph.D - Information Engineering and Computer Science} \textit{\small{ - University of Trento, Italy}} \hfill \textit{Jan 2025}\\[2pt]
\noindent\textit{\small Thesis: Embedded Real-time Musical Pattern Detection for Smart Musical Instruments}\\[3pt]
\noindent\textbf{M.Sc - Telecommunication and Electronic Engineering} \textit{\small{ - Sheffield Hallam University, UK}} \hfill \textit{Aug 2018}\\[3pt]
\noindent\textbf{B.Eng (Hons) - Electronic Engineering} \textit{\small{ - Sheffield Hallam University, UK}} \hfill \textit{Mar 2014}

\section*{Research and Work Experience}

\begin{cvrole}{Independent Researcher}{Self-directed}{Toronto, Canada}{Jan 2026}{Present}
    \item \textbf{VarianceEngine}: generative model producing expressive musical variations from a single ground-truth recording; compared 5 architectures (mel-diffusion, VQ-VAE + transformer prior, retrieval baseline, symbolic transcription + transformer, foundation-model fine-tune) before choosing LoRA fine-tuning of MusicGen-medium with a custom audio-to-audio conditioning head, on dual-GPU (RTX 4090 + 3090). In progress.
    \item \textbf{speakfrench}: benchmarked 6 pronunciation-scoring methods (DTW/cosine-MFCC, Gaussian log-likelihood/KL, Wav2Vec2) at scale - 100,000 scored trials, 10,000 sentences, 5 voices - via ROC-AUC, EER, Cohen's~d.
    \item \textbf{melodiq}: audio-feature ETL pipeline (Airflow, S3-to-Snowflake, Dockerized, RDS/Snowflake infra-as-code). \textbf{poetry2music} (in progress): multilingual annotated poetry corpus, phonetic/prosodic feature-annotation pipeline.
    \item Two papers accepted this period (Smart Drums, JAES; MIRaaS, IEEE IS\textsuperscript{2} - see Publications); maintain \textbf{Nebula}, a C++17 audio-feature library with ARM/Raspberry Pi latency benchmarks.
\end{cvrole}

\needspace{4\baselineskip}
\begin{cvrole}{Postdoctoral Researcher}{University of Trento}{Trento, Italy}{Jan 2025}{Dec 2025}
    \item Owned the science and the full ML pipeline for a novel real-time system on streaming audio and IMU/gesture time series: acquisition, DSP and feature engineering, sequence-model training and evaluation, edge deployment with monitoring; EU Horizon EIC Pathfinder project MUSMET.
    \item Designed frozen-protocol benchmark suites and ablations (RNN vs DTW; audio vs MIDI; pattern-set sizes 1/3/10) to separate real accuracy gains from latency and CPU noise, with a written record of negative results; published system F1 0.76 at 14 ms on a Raspberry Pi 4, 31.4\% CPU, 74$\times$ faster than the 1038 ms DTW baseline (JAES 2026).
    \item Prototyped VAE- and diffusion-based synthetic time-series generation to expand limited-label training data, and a rule-based variation generator used across two published studies.
    \item Built a real-time multimodal synchronization pipeline aligning audio, BCI/EEG (g.tec) and mixed-reality head/hand tracking across four musicians simultaneously, with SAE Institute Barcelona; ran a user study across two live concerts (20 audience, 6 performers, p $<$ 0.05).
\end{cvrole}

\needspace{4\baselineskip}
\begin{cvrole}{Visiting Researcher}{McGill University (IDMIL / CIRMMT)}{Montreal, Canada}{May 2024}{Aug 2024}
    \item Fused inertial (IMU) and audio time series through a hierarchical finite-state machine, accumulating repeated observations of the same event to reach F1 0.78 across a 500-event corpus (IEEE IS\textsuperscript{2} 2025).
    \item Profiled compute and power for real-time feasibility on edge-class hardware.
\end{cvrole}

\needspace{4\baselineskip}
\begin{cvrole}{Visiting Researcher}{University of Visual and Performing Arts}{Colombo, Sri Lanka}{Jan 2024}{Mar 2024}
    \item Designed and ran a human-centred evaluation of ML-driven interactive systems, combining performance metrics with structured interviews to produce usability, interpretability and adoption recommendations (IEEE I3DA 2025).
\end{cvrole}

\needspace{4\baselineskip}
\begin{cvrole}{Technical Consultant}{Forestpin (Pvt) Ltd}{Colombo, Sri Lanka}{Aug 2020}{Feb 2021}
    \item Designed and deployed ML and statistical pipelines for financial forensics, compliance monitoring and risk detection on large, messy structured enterprise datasets.
    \item Built SQL and Python backend anomaly-scoring services running in production triage; translated compliance-stakeholder requirements into technical specifications.
\end{cvrole}

\needspace{4\baselineskip}
\begin{cvrole}{Research Engineer}{MAS Holdings (Pvt) Ltd}{Colombo, Sri Lanka}{Jan 2015}{Jul 2020}
    \item Framed and shipped end-to-end computer-vision and IoT systems for manufacturing at scale, moving research prototypes to deployed production workflows with CI/CD and code review; operational efficiency improved by up to 300\%.
    \item Built high-frequency IoT ETL from camera and sensor streams across factories in South Asia, North America and Africa, feeding live production monitoring and analytics.
    \item Delivered automated care-label QC (OpenCV template and feature matching, conveyor with PLC reject) across three hardware generations; inspection time cut 99.5\% with a human-in-the-loop path for borderline cases. Loom-side fabric inspection with a microscopic camera array cut operator headcount 50\%.
    \item Ran group-wide computer-vision training, mentored engineers, and delivered a COVID-19 remote-vitals CV device deployed across hospital wards (GMOA recognition).
\end{cvrole}

\section*{Publications \normalfont{\footnotesize (Google Scholar: 27 citations, h-index 3)}}
\begin{itemize}[leftmargin=1.2em, itemsep=-0.25em]
    \item Silva, N. and Turchet, L. \textit{Real-Time Audio Pattern Detection for Smart Musical Instruments}. Journal of the Audio Engineering Society, Vol.~74, No.~3, Mar.~2026.
    \item Silva, N. and Turchet, L. \textit{Smart Drums: Comparing Audio and MIDI in Embedded Real-Time Drum Pattern Recognition}. Journal of the Audio Engineering Society, 2026 (in press).
    \item Silva, N. and Turchet, L. \textit{Music Information Retrieval as a Service: Latency Characterization of an Edge-to-Server Architecture for Near-Real-Time Audio Analysis}. IEEE International Symposium on the Internet of Sounds, 2026 (accepted).
    \item Silva, N., Boem, A. and Turchet, L. \textit{Interactive IoMusT-Based Concerts: Real-Time Pattern Recognition and Audience Experience}. IEEE International Conference on Immersive and 3D Audio, Sep.~2025.
    \item Silva, N. and Turchet, L. \textit{Real-Time Pattern Recognition of Symbolic Monophonic Music}. 19th International Audio Mostly Conference (ACM), Sep.~2024.
    \item Silva, N. and Turchet, L. \textit{A Structural Similarity Index Based Method to Detect Symbolic Monophonic Patterns in Real-Time}. 25th International Conference on Digital Audio Effects (DAFx), Sep.~2022.
    \item Silva, N., Weeraddana, P.C. and Fischione, C. \textit{On Musical Onset Detection via the S-Transform}. 52nd Asilomar Conference on Signals, Systems, and Computers, Oct.~2018.
    \item Silva, N. \textit{Embedded Real-time Musical Pattern Detection for Smart Musical Instruments}. PhD Thesis, University of Trento, Jan.~2025.
\end{itemize}

\section*{Datasets \normalfont{\footnotesize (open access, Zenodo)}}
Four open-access benchmark datasets, 9,800+ recordings and 70+ musicians, each with a frozen evaluation protocol: \textbf{DoMP} (4,392 monophonic MIDI recordings, 40 musicians), \textbf{DoPP} (2,276 polyphonic audio recordings, 44.1~kHz/24-bit, 20 musicians), \textbf{DoDP} (994 drum-pattern MIDI recordings, 10 drummers), \textbf{DoDP2} (2,177 drum-pattern MIDI recordings, 10 drummers, artist-level style templates). DOIs: \href{https://doi.org/10.5281/zenodo.10818617}{\textit{10818617}}, \href{https://doi.org/10.5281/zenodo.14497998}{\textit{14497998}}, \href{https://doi.org/10.5281/zenodo.14497974}{\textit{14497974}}, \href{https://doi.org/10.5281/zenodo.18395007}{\textit{18395007}}.

\enlargethispage{3\baselineskip}
\section*{Service, Review \& Supervision}
\begin{itemize}[leftmargin=1.2em, itemsep=-0.25em]
    \item \textbf{Reviewer:} IEEE Access; Journal of the National Science Foundation, Sri Lanka; Colloquio di Informatica Musicale. \textbf{Program Committee:} IEEE International Symposium on the Internet of Sounds; International Conference on Immersive and 3D Audio.
    \item \textbf{BSc Thesis Supervision, University of Trento:} Leonardo Collizoli - MIDI plugin for real-time musical pattern recognition (Aug 2023); Dennis Battisti - plugin and GUI for real-time musical pattern detection (Sep 2026).
\end{itemize}

\section*{Highlights and Recognition}
\textbf{MIDI Innovation Awards:} Finalist (2023), Prototypes \& Non-Commercial Software - ``Hot Licks'' real-time pattern detection for smart instruments. \textbf{GMOA Recognition (Sri Lanka):} COVID-19 remote-vitals monitoring device deployed across hospital wards. \textbf{Open source:} LiveLaTeX, a published VS Code Marketplace extension; side projects built with AI-assisted workflows including Claude Code.

\end{document}
